\section{Discussion}

The central distinction is between a log that records application activity and an evidence object that supports scientific reconstruction. The latter must preserve not only what happened but which prespecified rules made the event eligible, how time was measured, why the system adapted, and why an observation was treated as recovered, censored, invalid, or missing.

\subsection{What conformance can falsify}

Three evidence layers must remain separate. Structural conformance tests whether an event satisfies the declared field and type constraints. Semantic conformance tests cross-field invariants, including terminal-state exclusivity, reference identity and timing eligibility. Deterministic replay tests whether a pinned controller and configuration reproduce expected decisions from an ordered event sequence. A passing result in one layer does not establish a passing result in either of the others.

A scientifically invalid measurement context may still be represented by a structurally valid event if its invalidity is explicit. Conversely, a syntactically well-formed event with mutually inconsistent recovery and censoring states must be rejected. Negative cases should therefore distinguish \emph{correct representation of invalid data} from \emph{acceptance of an invalid representation}. The declared corpus includes both kinds of case, but its execution and diagnostic specificity must be demonstrated at the cited revision before reporting results.

\subsection{Timing and transport threats}

A monotonic within-session clock prevents some wall-clock artifacts but does not, by itself, establish timestamp precision, inter-device synchronization or continuity across a device restart. Prospective cross-device use requires a documented clock-domain mapping, reset policy, sampling-resolution bound and treatment of missing or reordered records. Queue acknowledgement and duplicate-delivery tests address transport behavior; they cannot substitute for recovery-event validation. These threats are requirements for further tests, not measured performance claims of the current implementation.

\subsection{External falsification criteria}

A cross-application portability claim would be falsified if an independently implemented serious game cannot encode its challenge eligibility, regulation action, recovery reference, observation window, and terminal-state semantics without either discarding information or silently reinterpreting an existing field. A deterministic-replay claim would be falsified if two executions at the same pinned controller revision and configuration yield different decisions for an identical ordered input sequence, after explicitly accounting for documented nondeterministic inputs. A transport-integrity claim would be falsified if duplicate delivery, reordering, or acknowledgement loss changes the reconstructed scientific outcome without a detectable provenance or validation error. These are proposed independent evaluation criteria, not results of tests already conducted.

An external evaluation should prespecify the independent game's event mapping and exclusion rules before testing, retain every failed mapping and replay trace, and report each validation layer separately. Internal fixtures alone cannot establish portability, and a successful mapping cannot establish construct equivalence or clinical validity. Any extension of the schema required to pass the external evaluation must receive a new version and a fresh, independently assessed test set rather than being counted as confirmation of the original frozen contract.

\paragraph{Independent portability audit design.}
To avoid selecting only applications that already fit the contract, an external audit should register candidate games and inclusion criteria before inspecting mapping outcomes. For each eligible application, two independent implementers should document the source event vocabulary, mapping to schema fields, any information loss, and disagreements resolved against the original application specification. The audit record should distinguish a missing source concept from a mapping failure and from a schema validation failure. Report the denominator of eligible applications, the number mapped without schema changes, the number requiring extensions, and all excluded cases with reasons. No external application sample, inter-rater agreement estimate, or portability rate is available in the present work; these remain prospective tests rather than established results.

\subsection{Outcome-preserving replay equivalence}
For reproducibility, a replay comparator should evaluate scientific outcomes rather than byte-for-byte equality of transport envelopes. Before execution, it must specify an event identity rule, the ordered decision inputs, permitted volatile transport metadata, and a canonical representation of each terminal outcome. Two traces are equivalent only if they preserve the challenge and recovery references, eligibility decisions, controller-state transitions, observation-window boundaries, and the distinction among observed recovery, right censoring, invalidity, and missingness. Duplicate transport deliveries may be collapsed only when they share a verified event identity and do not alter decision inputs; distinct events with equal payloads must not be collapsed. An observed recovery must never be equated with a censored outcome merely because their reported durations match. This is a prespecified comparison criterion for future independent replay, not a claim that such a comparison has passed.

\subsection{Limits of transfer and inference}

The schema does not make heterogeneous interventions scientifically equivalent. A common event structure cannot rescue incompatible challenge definitions, different constructs, weak measurement models or unmeasured confounding. Generalization beyond ARC-ADHD requires a study-specific mapping of challenge eligibility, recovery reference, stability criterion, observation horizon and censoring rule, followed by independent testing. In particular, an event that can be replayed does not show that the corresponding recovery endpoint measures self-regulation or predicts functional improvement.

A reusable challenge--regulation--recovery representation could improve algorithm comparison, provenance across application versions, missingness interpretation and reconstruction of adaptive-controller decisions. These are methodological possibilities. Claims of interoperability, measurement validity or intervention efficacy require evidence beyond the present schema and software fixtures.
